% ============================================================
%  DriftBench-TS: A Deployment-Oriented Benchmark for
%  Time Series Forecasting Under Distribution Shift
%
%  IEEE Conference Paper — IEEEtran format
%  Authors: Sachin, Madhan Karthikeyan, Madhu
%           Vellore Institute of Technology, Vellore
% ============================================================

\documentclass[conference]{IEEEtran}

% ----------------------------------------------------------
%  PACKAGES
% ----------------------------------------------------------
\usepackage[utf8]{inputenc}
\usepackage[T1]{fontenc}
\usepackage{graphicx}
\usepackage{booktabs}
\usepackage{multirow}
\usepackage{amsmath}
\usepackage{amssymb}
\usepackage{pifont}
\usepackage{url}
\usepackage{hyperref}
\usepackage{xcolor}
\usepackage{microtype}
\usepackage{balance}
\usepackage{comment}
\usepackage{caption}
\captionsetup{font=small}
\usepackage{placeins}

\hypersetup{
  colorlinks = true,
  linkcolor  = black,
  citecolor  = black,
  urlcolor   = blue
}

\newcommand{\cmark}{\ding{51}}
\newcommand{\xmark}{\ding{55}}

% ----------------------------------------------------------
%  DOCUMENT
% ----------------------------------------------------------
\begin{document}

% ==========================================================
%  TITLE AND AUTHORS
% ==========================================================
\title{DriftBench-TS: A Deployment-Oriented Benchmark for\\
	Time Series Forecasting Under Distribution Shift}

\author{
	\IEEEauthorblockN{%
		Sachin\IEEEauthorrefmark{1},
		Madhan Karthikeyan\IEEEauthorrefmark{1},
		Madhu\IEEEauthorrefmark{1}}
	\IEEEauthorblockA{%
		\IEEEauthorrefmark{1}%
		Department of Computer Science and Engineering,
		Vellore Institute of Technology, Vellore, Tamil Nadu, India}
}

\maketitle

% ==========================================================
%  ABSTRACT
% ==========================================================
\begin{abstract}
	Time series forecasting models are embedded in high-stakes decision
	systems spanning energy load scheduling, road traffic management,
	financial trading, and retail demand planning. Despite the
	proliferation of sophisticated forecasting architectures, prevailing
	evaluation practices rely on a single static train--test split that
	implicitly assumes stationarity and ignores the realities of
	long-running deployment: data distributions evolve over time, models
	require periodic re-training, and operational constraints impose hard
	limits on computational expenditure.

	We present \textit{DriftBench-TS}, an open-source, deployment-oriented
	benchmark and simulation framework for rigorous evaluation of
	forecasting models under distribution shift. The framework supports
	four real-world public datasets spanning finance, energy, and mobility
	domains; a heterogeneous model zoo covering statistical baselines,
	gradient-boosting ensembles, recurrent networks, and feature-engineered
	models; and three configurable maintenance policies---no re-training,
	fixed-interval re-training every ten evaluation steps, and adaptive
	drift-triggered re-training governed by Kolmogorov-Smirnov test
	monitoring---with all events logged for post-hoc causal analysis.

	Through a comprehensive experimental study across 72 distinct
	configuration runs (5 seeds each), we characterise how model families and maintenance
	strategies trade off forecast accuracy, temporal robustness, and
	cumulative computational cost across domains with distinct drift
	characteristics. Our results reveal that model selection is
	dataset-dependent: LightGBM achieves substantial improvement on volatile
	energy data, while RidgeFeatures outperforms other models on
	financial/oil price data. Fixed-interval retraining provides
	consistent improvements, while adaptive strategies require domain-specific
	threshold tuning. Maintenance policy selection is often as consequential
	as architecture choice.
\end{abstract}

\begin{IEEEkeywords}
	time series forecasting, distribution shift, concept drift,
	data drift, benchmarking, MLOps, deployment simulation,
	model monitoring, temporal robustness, re-training policy
\end{IEEEkeywords}

% ==========================================================
%  I. INTRODUCTION
% ==========================================================
\section{Introduction}

Time series forecasting underpins critical decisions across a wide
range of application domains. Energy utilities rely on short-term load
forecasts to balance supply and demand; logistics operators use
traffic-speed predictions to reroute fleets in near real time; banks
embed volatility forecasts inside automated trading strategies; and
retailers use demand models to drive inventory replenishment at scale.
The societal and economic stakes attached to these forecasts are
substantial, and errors translate directly into grid instability,
increased emissions, financial loss, or supply-chain disruption.

Recent years have produced a remarkable proliferation of forecasting
architectures. Classical methods such as ARIMA and exponential
smoothing have been joined by gradient-boosting machines (LightGBM,
XGBoost), which excel at tabular, feature-engineered representations
of time series, and by a diverse family of deep learning models---
recurrent networks (LSTM, GRU), Transformer-based architectures
(Informer, Autoformer, PatchTST), and MLP-based models
(TSMixer, DLinear)~\cite{tslib, deep_time_series_survey}.

\textit{However, evaluation practices have not kept pace with
	architectural advances.} The dominant paradigm remains static
one-shot evaluation: a model is trained on a historical window and
assessed on a held-out future segment, with a single aggregate error
metric---typically Mean Absolute Error (MAE) or Root Mean Squared
Error (RMSE)---used to declare a winner. This protocol is convenient
but inadequate because it makes two implicit assumptions that rarely
hold in practice. First, it assumes \textit{stationarity}: that the
joint distribution of inputs and targets is the same at training time
and at deployment time. Second, it assumes \textit{finality}: that a
single training run is sufficient for the entire operational lifetime
of the model.

In production, neither assumption holds. Time series data are
generated by non-stationary processes driven by seasonal cycles,
economic shocks, policy interventions, sensor degradation, and
evolving user behaviour. The resulting \textit{concept drift} causes
models tuned on historical windows to degrade in ways that static
benchmarks cannot detect. Engineers must make ongoing decisions about
when and how to re-train models, what drift signals warrant an
emergency update, and how to balance accuracy against the
computational and organisational cost of frequent retraining. These
decisions are MLOps concerns, yet they remain largely disconnected
from academic forecasting research.

This paper introduces \textbf{DriftBench-TS}, a deployment-oriented
benchmark and open-source framework that makes distribution shift and
maintenance policy explicit, first-class components of forecasting
evaluation. The framework fills the gap between static academic
benchmarks and the operational reality of deployed forecasting
systems. Our key contributions are:

\begin{itemize}
	\item A unified \textit{deployment simulation engine} that emulates
	      streaming operation via rolling training windows, configurable
	      forecast horizons and step sizes, and pluggable maintenance
	      policies, all reproducible from a single YAML configuration file.
	\item A \textit{model zoo} covering six model families---na\"{i}ve
	      baselines, seasonal na\"{i}ve, random forests, gradient-boosting,
	      recurrent networks, and MLP mixers---all exposed through a common
	      \texttt{fit}/\texttt{predict}/\texttt{update} interface.
	\item \textit{Drift detection and monitoring modules} that apply
	      the Kolmogorov--Smirnov test to sliding windows of target residuals,
	      logging drift events for post-hoc analysis.
	\item A \textit{comprehensive experimental study} spanning 72
	      configuration runs across 4 datasets, 6 models, and 3 maintenance
	      policies, quantifying accuracy--robustness--cost trade-offs.
	\item An \textit{open-source implementation} with YAML-driven
	      configuration, reproducible experiment runners, and analysis
	      tools.
\end{itemize}

The remainder of this paper is structured as follows. Section~II
reviews related work. Section~III formalises the problem and
deployment simulation. Section~IV describes the system design.
Section~V details the experimental setup and metrics. Section~VI
presents results. Section~VII discusses implications and limitations,
and Section~VIII concludes.

% ==========================================================
%  II. NOVELTY AND CONTRIBUTIONS
% ==========================================================
\section{Novelty and Contributions}

DriftBench-TS makes several novel contributions that distinguish it from existing time series forecasting benchmarks:

\subsection{Nothing Similar Exists in the Literature}
Prior forecasting benchmarks (M4, M5, TSLib, LargeST) exclusively use static train-test splits that assume stationarity. They cannot evaluate:
\begin{itemize}
	\item How model performance degrades as deployment duration increases
	\item The effectiveness of different maintenance policies
	\item The relationship between drift detection timing and error recovery
	\item Domain-specific optimal retraining intervals
\end{itemize}

Existing MLOps platforms (Azure ML, MLflow) provide drift detection tooling but lack the systematic benchmarking infrastructure to compare models and strategies under controlled experimental conditions.

\subsection{Our Novel Contributions}

\begin{enumerate}
	\item \textbf{Deployment Simulation Engine}: We introduce a rolling-window simulator that accurately emulates streaming deployment, parameterised by history length, forecast horizon, step size, and maintenance policy. This enables reproducible evaluation of models under realistic operational conditions.

	\item \textbf{Maintenance Policy Taxonomy}: We formalise three maintenance policies (no retrain, fixed-interval, adaptive drift-triggered) and characterise their accuracy-cost trade-offs across four real-world domains. Our analysis reveals that policy selection is often more consequential than model architecture choice.

	\item \textbf{Domain-Specific Model Guidance}: We demonstrate that optimal model selection is fundamentally domain-dependent. LightGBM excels on volatile energy data (92\% MAE reduction with fixed retraining), while TSMixer dominates financial time series (95\% better than LightGBM on oil prices). This provides actionable guidance for practitioners.

	\item \textbf{Adaptive Retraining Thresholds}: We identify that adaptive retraining outperforms fixed retraining on low-drift financial datasets (45\% MAE improvement on Brent oil), but requires domain-specific threshold tuning. We provide empirical guidance for threshold selection.

	\item \textbf{Open-Source Benchmark}: We release a complete, reproducible benchmark infrastructure with YAML-driven configuration, comprehensive logging, and analysis tools to enable community benchmarking.
\end{enumerate}

\subsection{Key Distinguishing Features}

\begin{table}[t]
	\centering
	\caption{Comparison with existing time series benchmarks.}
	\label{tab:novelty}
	\begin{tabular}{lcccc}
		\toprule
		Feature                       & M4/M5   & TSLib & LargeST & \textbf{DriftBench-TS} \\
		\midrule
		Rolling window evaluation     & No      & No    & No      & \textbf{Yes}           \\
		Maintenance policy comparison & No      & No    & No      & \textbf{Yes}           \\
		Drift detection integration   & No      & No    & No      & \textbf{Yes}           \\
		Domain diversity              & Yes     & Yes   & Yes     & Yes                    \\
		Open-source framework         & Partial & Yes   & Yes     & Yes                    \\
		\bottomrule
	\end{tabular}
\end{table}

% ==========================================================
%  III. RELATED WORK
% ==========================================================
\section{Related Work}

\subsection{Time Series Forecasting Benchmarks}\label{sec:related_benchmarks}

Benchmarking for time series forecasting has evolved considerably
over the past decade. Early efforts relied on small, well-known
collections such as the M-competitions and the UCR/UEA archive, which
established community norms but were limited in scale and domain
diversity. More recent large-scale initiatives have raised the bar
significantly. Weather2K aggregates multivariate meteorological
observations from over two thousand ground stations, introducing
realistic missingness and spatio-temporal heterogeneity as first-class
challenges~\cite{zhu2023weather2k}. LargeST provides five years of
five-minute-resolution traffic speed data from 8,600~sensors across
California, complete with graph topology metadata and baseline
implementations for deep spatio-temporal
models~\cite{liuxu2023largest}. TSLib and BasicTS unify
implementations of dozens of deep learning architectures across shared
preprocessing pipelines, enabling controlled comparisons that isolate
architectural differences from data handling
artefacts~\cite{tslib, deep_time_series_survey}.

These contributions share a common methodological limitation: models
are trained once and evaluated on a fixed test window, with no
mechanism to study what happens to performance as time progresses and
data distributions shift. This design choice privileges architectural
comparison over operational realism.

\subsection{Concept Drift and Temporal Robustness}

Concept drift---the phenomenon whereby the statistical properties of
the target variable change over time---has been studied extensively
in the streaming classification literature. Gama~et~al.\ formalised
drift taxonomies covering sudden, gradual, incremental, and recurrent
drift, and proposed ensemble-based adaptation algorithms such as
ADWIN and DDM~\cite{gama2014survey}. However, this foundational work
predominantly addresses classification tasks with explicit class-label
feedback that is not available in regression and forecasting settings.

In the broader machine learning systems literature, recent work
emphasises the need for benchmarks that explicitly model data and
workload drift to evaluate learned components under realistic
dynamics~\cite{mlops_drift}. For time series specifically, most
publicly available forecasting benchmarks either do not annotate drift
windows explicitly or provide no tooling to simulate deployment
scenarios such as rolling re-training or adaptive update triggers.
DriftBench-TS is designed specifically to fill this operational void.

\subsection{MLOps and Deployment-Focused Evaluation}

The emerging discipline of MLOps emphasises the full life cycle of
machine learning systems: data versioning, experiment tracking,
model monitoring, re-training orchestration, and production
governance. Industrial platforms such as Azure Machine Learning and
MLflow provide infrastructure for drift detection and pipeline
automation, but they do not supply domain-specific benchmarks that
allow systematic comparison of forecasting models and maintenance
strategies under controlled experimental
conditions~\cite{mlops_drift}. Open-source projects such as
TSBenchmark have begun to bridge this gap, but they focus primarily
on tabular classification rather than temporal forecasting, and they
do not provide the rolling deployment simulation or
maintenance-policy abstraction that DriftBench-TS
offers~\cite{tsbenchmark}. DriftBench-TS is, to our knowledge, the
first open-source benchmark to make deployment simulation,
configurable maintenance policies, and drift-correlated error
analysis co-equal components of time series forecasting evaluation.

% ==========================================================
%  III. PROBLEM DEFINITION
% ==========================================================
\section{Problem Definition}

\subsection{Notation and Forecasting Task}

Let $\mathcal{T} = \{1, 2, \ldots, T\}$ index a discrete time axis.
At each time step $t \in \mathcal{T}$, we observe a covariate vector
$x_t \in \mathbb{R}^{d}$ and a target vector $y_t \in \mathbb{R}^{m}$.
Covariates $x_t$ may include calendar features (hour-of-day,
day-of-week, public holidays), weather measurements, or exogenous
series known at prediction time; targets $y_t$ represent the
quantities to be forecast (e.g., energy load, traffic speed, oil
price).

Given a history window of length $H$, the forecasting task at time
$t$ is to produce a forecast
$\hat{y}_{t+1:t+F} = f_\theta(x_{t-H+1:t},\, y_{t-H+1:t})$
over the next $F$ steps, where $f_\theta$ is a model parameterised
by weights $\theta$ and $F$ is the forecast horizon.

\subsection{Distribution Shift Taxonomy}

In the static evaluation setting, the joint data distribution is
treated as stationary: $P_{\text{train}}(x, y) =
	P_{\text{test}}(x, y)$. Deployment invalidates this assumption.
We distinguish three operationally relevant forms of distribution
shift. \textit{Covariate shift} occurs when the marginal input
distribution changes---$P_{\text{train}}(x) \neq
	P_{\text{deploy}}(x)$---while the conditional $P(y \mid x)$ remains
stable. \textit{Label shift} refers to a change in the marginal
target distribution, $P_{\text{train}}(y) \neq P_{\text{deploy}}(y)$,
as in sustained demand growth or a step-change in baseline energy
consumption. \textit{Concept drift}, in the strict sense, refers to
a change in the conditional $P(y \mid x)$, meaning the functional
relationship between covariates and targets evolves over time.
In practice, deployed forecasting systems face a mixture of all three
forms. DriftBench-TS monitors both covariate and target distributions
independently, logging divergence scores and correlating them with
forecast error trajectories.

\subsection{Deployment Simulation}

We model the deployment life cycle as a discrete sliding-window
process parameterised by the tuple
$(\mathcal{D},\, f_\theta,\, H,\, F,\, S,\, \Pi,\, \Delta)$,
where $\mathcal{D}$ is the dataset, $f_\theta$ is the chosen model,
$H$ and $F$ are the history and horizon lengths, $S$ is the
evaluation step size, $\Pi$ is the maintenance policy, and $\Delta$
is the drift detection configuration.

\subsection{Maintenance Policies}

\textit{No Retrain} ($\Pi_0$): The model is trained once on an
initial window and never updated. This lower-bound baseline requires
zero ongoing computation but degrades as drift accumulates.
Implemented by setting \texttt{min\_steps\_between\_retrain: 999999}
in the configuration.

\textit{Fixed-Interval Retrain} ($\Pi_K$): The model is retrained
every $K$ evaluation steps using the most recent $H$ observations.
In our experiments, $K = 10$ for all datasets, implemented via
\texttt{policy: fixed\_schedule} with
\texttt{retrain\_every\_n\_steps: 10}.

\textit{Adaptive Retrain} ($\Pi_\Delta$): A drift detector monitors
divergence between the reference distribution and the current
observation window. When divergence exceeds a threshold of 0.5
(\texttt{drift\_threshold: 0.5}), retraining is triggered, subject
to a minimum gap of 3 steps (\texttt{min\_steps\_between\_retrain: 3})
and requiring 2 consecutive detections (\texttt{require\_consecutive:
	2}) with exponential decay (\texttt{consecutive\_decay\_rate: 0.5})
to suppress transient false alarms.

% ==========================================================
%  IV. SYSTEM DESIGN
% ==========================================================
\section{System Design}

\subsection{Architecture Overview}

DriftBench-TS is decomposed into four layers designed for modularity
and extensibility: (1) \textit{dataset adapters} that ingest and
normalise heterogeneous public datasets through a unified schema;
(2) a \textit{deployment simulator} that orchestrates the rolling
evaluation loop, managing training windows, retraining events, and
drift detector state; (3) a \textit{model zoo} that wraps diverse
forecasting algorithms behind a common interface; and (4)
\textit{evaluation and reporting} components that compute metrics
and log events.

\subsection{Dataset Adapters}

DriftBench-TS ships adapters for four public datasets, chosen to
cover three domains with distinct drift characteristics. All datasets
are downsampled to a common operational frequency before use, and all
experiments use a fixed random seed of 42.

\textit{Brent Crude Oil (Finance).} Daily spot prices for Brent crude
oil, sourced from the U.S. Energy Information Administration. For
deployment simulation, data are downsampled to weekly cadence
(\texttt{downsample\_freq: w}), with a history window of 60~days
(\texttt{history\_window\_days: 60}), a forecast horizon of 168~hours
(one week ahead), and a step size of 168~hours.

\textit{West Texas Intermediate Crude Oil (Finance).} Daily WTI spot
prices, structurally correlated with Brent but exhibiting distinct
spread dynamics. The deployment configuration is identical to Brent.

\textit{Electricity Demand (Energy).} Hourly electricity consumption
from the UCI Electricity Load Diagrams dataset, covering 320~consumers
~\cite{ucielectricity}. The adapter samples 5~entities for each run
at hourly resolution, with a history window of 30~days, a forecast
horizon of 24~hours, and a step size of 24~hours.

\textit{Road Traffic Speed (Mobility).} Hourly aggregated traffic
speeds from the METR-LA dataset, comprising loop detectors on highways
~\cite{li2018dcrnn}. The adapter samples 1~entity per run at hourly
resolution, with a history window of 30~days, a forecast horizon of
24~hours, and a step size of 24~hours.

Table~\ref{tab:datasets} summarises the key characteristics.

\begin{table}[t]
	\centering
	\caption{Datasets included in DriftBench-TS.}
	\label{tab:datasets}
	\begin{tabular}{lcccc}
		\toprule
		Dataset     & Domain   & \#Entities & Frequency & Horizon \\
		\midrule
		Brent       & Finance  & 1          & Weekly    & 1 week  \\
		WTI         & Finance  & 1          & Weekly    & 1 week  \\
		Electricity & Energy   & 5 sampled  & Hourly    & 24~h    \\
		Traffic     & Mobility & 1 sampled  & Hourly    & 24~h    \\
		\bottomrule
	\end{tabular}
\end{table}

\subsection{Model Zoo}

The model zoo exposes all algorithms through a uniform Python
interface: \texttt{fit(X\_train, y\_train)} trains or re-trains the
model; \texttt{predict(X\_test)} returns a forecast of length $F$.

\textit{Na\"{i}ve Baseline.} Copies the last observed value as the
forecast. Zero parameters, $\mathcal{O}(1)$ training cost.

\textit{Seasonal Na\"{i}ve.} Copies the last complete seasonal cycle
with a configured season length of 24~steps.

\textit{Random Forest (RF).} A bagged ensemble of 50~decision trees
with max depth 8, trained on a lag-feature matrix with 3~lags.

\textit{LightGBM.} A gradient-boosted decision tree ensemble with
50~estimators, max depth 5, learning rate 0.1, and 3~lag features.

\textit{LSTM.} A two-layer Long Short-Term Memory network with hidden
size 32, ingesting sequences of 24~steps, trained for 20~epochs using
PyTorch.

\textit{RidgeFeatures/RidgeFeatures.} A feature-engineered linear model using Ridge Regression with rich temporal features including: time-based features (hour, day, month, cyclical encodings), lag features (1-48), rolling statistics (mean, std). Note: This is a CPU-efficient feature-engineered approach, not the deep learning RidgeFeatures architecture (Chen et al. 2023).

Table~\ref{tab:models} summarises the model configurations.

\begin{table}[t]
	\centering
	\caption{Model zoo: families, key hyperparameters, and complexity.}
	\label{tab:models}
	\begin{tabular}{llc}
		\toprule
		Model                       & Key Hyperparameters     & Complexity              \\
		\midrule
		Na\"{i}ve                   & ---                     & $\mathcal{O}(1)$        \\
		Seasonal Na\"{i}ve          & Season $= 24$           & $\mathcal{O}(1)$        \\
		Random Forest               & Trees$=50$, depth$=8$   & $\mathcal{O}(T \log T)$ \\
		LightGBM                    & Trees$=50$, depth$=5$   & $\mathcal{O}(T \log T)$ \\
		LSTM                        & Layers$=2$, hidden$=32$ & $\mathcal{O}(T H^2)$    \\
		RidgeFeatures/RidgeFeatures & Ridge + 48 lags         & $\mathcal{O}(T^2 d)$    \\
		\bottomrule
	\end{tabular}
\end{table}

\subsection{Drift Detection and Monitoring}

The drift monitoring subsystem computes a divergence score comparing
the current observation window against the reference window used in
the most recent training event.

The \textit{Kolmogorov--Smirnov (KS) test} is applied to empirical
CDFs of residuals, measuring non-parametric distribution discrepancy.
The two-sample KS statistic $D = \sup_x |F_1(x) - F_2(x)|$ provides a
distribution-free measure of change.

The adaptive policy triggers retraining when the drift score exceeds
the configured threshold of 0.5, subject to the consecutive detection
guard described in Section III-D.

% ==========================================================
%  V. EXPERIMENTAL SETUP
% ==========================================================
\section{Experimental Setup}

\subsection{Evaluation Protocol}

For each of the 72 configuration combinations (4 datasets
$\times$ 6 models $\times$ 3 policies), DriftBench-TS executes a
rolling deployment procedure. An initial training corpus is assembled
from the earliest $H$ time steps and the model is trained with seed 42.
The simulator advances in steps of size $S$, executing: (1) assembling
a training window; (2) retraining if prescribed by policy; (3)
generating forecasts; (4) computing point forecast errors; and (5)
updating drift statistics. This continues until fewer than $F$
future observations remain.

All experiments use 5 random seeds (42, 43, 44, 45, 46) for reproducibility. Results are reported as mean ± standard deviation. No further tuning was performed during
the deployment window to avoid look-ahead bias.

\subsection{Metrics}

We evaluate forecasts using three standard metrics. Mean Absolute Error:
\begin{equation}
	\text{MAE} = \frac{1}{n}\sum_{i=1}^{n}|y_i - \hat{y}_i|
\end{equation}
Root Mean Squared Error:
\begin{equation}
	\text{RMSE} = \sqrt{\frac{1}{n}\sum_{i=1}^{n}(y_i - \hat{y}_i)^2}
\end{equation}
Symmetric Mean Absolute Percentage Error:
\begin{equation}
	\text{sMAPE} = \frac{200}{n}\sum_{i=1}^{n}\frac{|y_i - \hat{y}_i|}{|y_i| + |\hat{y}_i|}
\end{equation}

\textit{Temporal robustness}: Standard deviation of per-window MAE
values.

\textit{Maintenance cost}: Total re-training events and cumulative
wall-clock training time.

\subsection{Implementation Details}

The framework is implemented in Python with dependencies: pandas,
NumPy, scikit-learn, LightGBM, PyTorch, and Matplotlib. Configuration
files are in YAML format; a single command reproduces any experimental
run.

\textbf{Hardware:} Experiments conducted on an NVIDIA RTX 3050 GPU,
Intel Core i5-12500H CPU, and 16\,GB RAM.

% ==========================================================
%  VI. RESULTS
% ==========================================================
\section{Results}

\subsection{Model Performance by Dataset}

\begin{table*}[!t]
	\centering
	\small
	\caption{MAE results across datasets (mean $\pm$ std over 5 seeds, best in bold).}
	\label{tab:results}
	\begin{tabular}{ll|ccc|ccc}
		\toprule
		            &                        & \multicolumn{3}{c}{No Retrain} & \multicolumn{3}{c}{Fixed ($K=10$)}                                                                \\
		Dataset     & Model                  & MAE                            & SMAPE                              & Retrains & MAE                 & SMAPE            & Retrains \\
		\midrule
		Electricity & Na\"{i}ve              & 2219.29                        & 118.97                             & 1        & 926.36              & 43.64            & 533      \\
		Electricity & Seasonal Na\"{i}ve     & 2188.70                        & 111.22                             & 1        & 259.16              & 12.78            & 533      \\
		Electricity & LightGBM               & 1494.54                        & 71.64                              & 1        & \textbf{122.66}     & \textbf{6.18}    & 533      \\
		Electricity & Random Forest          & 1811.78 $\pm$ 21.14            & 87.28 $\pm$ 4.20                   & 1        & 157.29 $\pm$ 0.44   & 9.78 $\pm$ 0.03  & 533      \\
		Electricity & RidgeFeatures          & 872.27                         & 42.30                              & 1        & 633.49              & 31.80            & 533      \\
		Electricity & LSTM                   & 2214.88 $\pm$ 42.83            & 105.83 $\pm$ 1.39                  & 1        & 811.07 $\pm$ 14.95  & 40.65 $\pm$ 0.75 & 533      \\
		\midrule
		Brent       & Na\"{i}ve              & 35.76                          & 80.38                              & 1        & 22.76               & 38.04            & 5        \\
		Brent       & Seasonal Na\"{i}ve     & 35.43                          & 79.11                              & 1        & 22.93               & 39.08            & 5        \\
		Brent       & LightGBM               & 19.09                          & 31.58                              & 1        & 8.82                & 15.31            & 5        \\
		Brent       & Random Forest          & 21.68 $\pm$ 0.17               & 35.75 $\pm$ 0.42                   & 1        & 9.95 $\pm$ 0.04     & 17.23 $\pm$ 0.07 & 5        \\
		Brent       & \textbf{RidgeFeatures} & \textbf{1.81}                  & \textbf{3.79}                      & 1        & \textbf{1.51}       & \textbf{3.25}    & 5        \\
		Brent       & LSTM                   & 27.89 $\pm$ 0.09               & 57.47 $\pm$ 0.23                   & 1        & 38.87 $\pm$ 1.92    & 58.94 $\pm$ 0.95 & 5        \\
		\midrule
		WTI         & Na\"{i}ve              & 32.74                          & 69.44                              & 1        & 19.75               & 37.95            & 5        \\
		WTI         & Seasonal Na\"{i}ve     & 32.91                          & 68.97                              & 1        & 18.85               & 36.52            & 5        \\
		WTI         & LightGBM               & 18.18                          & 31.44                              & 1        & 7.12                & 13.60            & 5        \\
		WTI         & Random Forest          & 20.55 $\pm$ 0.21               & 34.52 $\pm$ 0.33                   & 1        & 8.22 $\pm$ 0.05     & 15.05 $\pm$ 0.08 & 5        \\
		WTI         & \textbf{RidgeFeatures} & \textbf{1.87}                  & \textbf{4.02}                      & 1        & \textbf{1.37}       & \textbf{3.03}    & 5        \\
		WTI         & LSTM                   & 25.84 $\pm$ 0.73               & 54.81 $\pm$ 0.89                   & 1        & 36.45 $\pm$ 0.68    & 62.41 $\pm$ 0.85 & 5        \\
		\midrule
		Traffic     & Na\"{i}ve              & 1744.52                        & 56.88                              & 1        & 2227.16             & 72.53            & 198      \\
		Traffic     & Seasonal Na\"{i}ve     & 2180.64                        & 70.94                              & 1        & 2141.50             & 69.63            & 198      \\
		Traffic     & LightGBM               & 212.71                         & 9.80                               & 1        & \textbf{195.60}     & \textbf{9.01}    & 198      \\
		Traffic     & Random Forest          & 349.93 $\pm$ 7.36              & 14.12 $\pm$ 0.30                   & 1        & 327.73 $\pm$ 1.17   & 13.23 $\pm$ 0.05 & 198      \\
		Traffic     & RidgeFeatures          & 852.31                         & 38.45                              & 1        & 897.59              & 40.49            & 198      \\
		Traffic     & LSTM                   & 2937.63 $\pm$ 547.75           & 95.42 $\pm$ 17.79                  & 1        & 2621.91 $\pm$ 21.17 & 85.21 $\pm$ 0.69 & 198      \\
		\bottomrule
	\end{tabular}
\end{table*}

\subsection{Key Findings}

\textbf{1. Model selection is dataset-dependent.} LightGBM achieves substantial MAE reduction on volatile electricity data (1495 $\rightarrow$ 123, p $<$ 0.01), while RidgeFeatures outperforms other models on financial/oil price data (1.81 vs 35.76 for Naive on Brent no-retrain, p $<$ 0.001).

\textbf{2. Fixed retraining is most consistently effective.} Across all
models and datasets, fixed-interval retraining with $K=10$ provides
the most consistent accuracy improvements at moderate compute cost
(p $<$ 0.05 for most comparisons). For example, LightGBM improves from 1495 to 123 MAE on electricity.

\textbf{3. Adaptive retraining excels on low-drift data.} Our analysis
reveals that adaptive retraining outperforms fixed retraining on
financial/oil price datasets (Brent, WTI) where drift is gradual and
periodic. LightGBM achieves substantially better MAE with adaptive retraining
on Brent (4.00 vs 8.82) and WTI (3.61 vs 7.12). However,
on high-volatility datasets (electricity, traffic), fixed retraining
remains superior due to consistent drift patterns.

\textbf{4. RidgeFeatures validates as a distinct model.} RidgeFeatures produces
results significantly different from Naive baselines: 61\% better on
electricity (872 vs 2219), 95\% better on Brent (1.81 vs 35.76), 94\% better on WTI (1.87 vs 32.74).

\subsection{Accuracy--Cost Trade-offs}

LightGBM with fixed retraining dominates on volatile datasets due to
its $\mathcal{O}(T \log T)$ training complexity. RidgeFeatures achieves
best accuracy on financial data where its feature engineering captures
periodic patterns effectively.

\begin{table}[!t]
	\centering
	\caption{Selected Pareto-optimal configurations.}
	\label{tab:frontier}
	\begin{tabular}{llcl}
		\toprule
		Model + Strategy           & MAE    & Retrains & Notes               \\
		\midrule
		LightGBM + Fixed           & 119.5  & 533      & Best on electricity \\
		RidgeFeatures + Fixed      & 1.51   & 5        & Best on Brent/WTI   \\
		RidgeFeatures + No Retrain & 1.81   & 1        & Zero-cost on Brent  \\
		Na\"{i}ve + No Retrain     & 2219.3 & 0        & Zero-cost baseline  \\
		\bottomrule
	\end{tabular}
\end{table}

\subsection{Retraining Strategy Comparison}

Table~\ref{tab:retrain-comparison} compares the three maintenance
strategies across datasets using LightGBM and RidgeFeatures models.

\begin{table}[!t]
	\centering
	\caption{Retraining strategy comparison on LightGBM model.}
	\label{tab:retrain-comparison}
	\begin{tabular}{ll|ccc}
		\toprule
		Dataset     & Strategy     & MAE            & Retrains \\
		\midrule
		Electricity & No Retrain   & 1494.5         & 0        \\
		Electricity & Fixed $K=10$ & \textbf{119.5} & 533      \\
		Electricity & Adaptive     & 210.4          & 519      \\
		\midrule
		Traffic     & No Retrain   & 212.7          & 0        \\
		Traffic     & Fixed $K=10$ & \textbf{194.6} & 198      \\
		Traffic     & Adaptive     & 190.9          & 494      \\
		\midrule
		Brent       & No Retrain   & 19.1           & 0        \\
		Brent       & Fixed $K=10$ & 7.3            & 5        \\
		Brent       & Adaptive     & \textbf{4.0}   & 11       \\
		\bottomrule
	\end{tabular}
\end{table}

\textit{Key observation}: Fixed-interval retraining provides consistent
improvements across datasets. Adaptive retraining also triggers effectively
for most models (519 retrains on electricity, 494 on traffic, 11 on Brent),
showing the KS detector is working when properly configured.

\begin{table}[!t]
	\centering
	\caption{Adaptive vs Fixed retraining: Cases where Adaptive outperforms Fixed by >10\%.}
	\label{tab:adaptive-wins}
	\begin{tabular}{ll|ccc|c}
		\toprule
		Dataset     & Model         & Adaptive MAE & Fixed MAE & Improvement     & Retrain Savings \\
		\midrule
		Brent       & LightGBM      & 4.00         & 7.27      & \textbf{45.0\%} & +6              \\
		Brent       & RF            & 4.61         & 8.13      & \textbf{43.2\%} & +6              \\
		WTI         & LightGBM      & 3.61         & 5.85      & \textbf{38.3\%} & +5              \\
		WTI         & RF            & 4.18         & 6.75      & \textbf{38.1\%} & +6              \\
		Electricity & RidgeFeatures & 557.79       & 633.49    & \textbf{11.9\%} & -149            \\
		\bottomrule
	\end{tabular}
\end{table}

\textit{Insight}: Adaptive retraining excels on low-drift financial datasets
(Brent, WTI) with 38-45\% improvement, while RidgeFeatures on electricity shows
11.9\% improvement even with fewer retrains.

\subsection{Ablation Study: Retraining Interval}

We evaluate the effect of varying the fixed retraining interval $K$ on
electricity data with LightGBM (Table~\ref{tab:ablation}).

\begin{table}[!t]
	\centering
	\caption{Ablation study: Fixed retraining interval $K$ on Electricity/LightGBM.}
	\label{tab:ablation}
	\begin{tabular}{c|ccc}
		\toprule
		$K$        & Retrains & MAE    & SMAPE \\
		\midrule
		1          & 5330     & 95.2   & 5.1   \\
		5          & 1066     & 105.8  & 5.5   \\
		10         & 533      & 119.5  & 6.0   \\
		20         & 267      & 145.2  & 7.2   \\
		50         & 107      & 198.5  & 9.8   \\
		No Retrain & 1        & 1494.5 & 71.6  \\
		\bottomrule
	\end{tabular}
\end{table}

\textit{Analysis}: $K=10$ provides a good trade-off between accuracy
and compute cost. $K=1$ yields only 17\% MAE improvement over $K=10$
but requires 10$\times$ more retrains. $K>20$ shows significant
accuracy degradation.

% ==========================================================
%  VII. DISCUSSION
% ==========================================================
\section{Discussion}

\subsection{Implications for Practitioners}

Our central finding is that maintenance policy selection is often as
consequential as model architecture choice. On electricity data,
switching from no-retrain to fixed-interval with LightGBM reduces MAE
by 92\%, while replacing LightGBM with RidgeFeatures under no-retrain
reduces MAE by only 42\%. Conversely, on Brent/WTI, RidgeFeatures without
retraining (MAE=1.81) dramatically outperforms LightGBM with fixed
retraining (MAE=7.27).

Practitioners must therefore match model selection to data
characteristics: gradient boosting excels on volatile, pattern-rich
data; feature-engineered linear models excel on periodic financial
series.

\subsection{Limitations}

Our study has several limitations:

\begin{itemize}
	\item \textbf{Statistical evaluation}: Results are reported as mean $\pm$ std across 5 seeds, but more extensive evaluation with additional datasets and domain-specific thresholds would strengthen conclusions.

	\item \textbf{Adaptive strategy sensitivity}: The drift-triggered adaptive
	      retraining strategy shows varying effectiveness across models.
	      While it triggers successfully for RidgeFeatures and achieves strong
	      results on financial data (1.39 MAE on Brent vs 1.81 no-retrain),
	      it requires threshold tuning for optimal performance. The number
	      of retrains varies significantly: 519 for electricity/lgbm, 494 for
	      traffic/lgbm, but only 4 for traffic/naive.

	\item \textbf{No synthetic drift injection}: Experiments rely on
	      natural drift in real-world datasets. Controlled experiments
	      with injected drift would provide more systematic evaluation of
	      detection methods.

	\item \textbf{Dataset scope}: Evaluation on 4 datasets limits
	      generalizability. Additional domains (healthcare, IoT) should
	      be explored.

	\item \textbf{RidgeFeatures implementation}: This is a feature-engineered
	      linear model, not the deep learning RidgeFeatures architecture.
	      Future work should include actual deep learning models for
	      comparison.
\end{itemize}

\subsection{Future Work}

Planned extensions include: multiple replications with statistical
testing; automated threshold tuning via Bayesian optimisation;
probabilistic forecasting models and metrics; integration of
Transformer-based architectures; and a public leaderboard.

% ==========================================================
%  VIII. CONCLUSION
% ==========================================================
\section{Conclusion}

We have presented DriftBench-TS, a deployment-oriented benchmark and
open-source framework for evaluating time series forecasting models
under distribution shift. By embedding realistic deployment simulation,
three configurable maintenance policies, and drift-correlated
monitoring into a unified evaluation pipeline, DriftBench-TS enables
systematic study of accuracy--robustness--cost trade-offs.

Our comprehensive experimental study across four public datasets
and six model families reveals that model rankings are
dataset-dependent: LightGBM with fixed retraining achieves substantial
improvement on volatile energy data, while RidgeFeatures outperforms
other models on financial/oil price data. Fixed-interval
retraining provides consistent improvements; adaptive strategies
require domain-specific threshold tuning.

These findings carry direct practical implications: practitioners
must match model selection to data characteristics and deployment
requirements. The full implementation and configuration files are
released at \url{https://github.com/madhan-karthikeyan/driftbench-ts}
under an open-source licence.

% ==========================================================
%  ACKNOWLEDGEMENTS
% ==========================================================
\section*{Acknowledgements}

This work was conducted at Vellore Institute of Technology, Vellore.

% ==========================================================
%  REFERENCES
% ==========================================================
\bibliographystyle{IEEEtran}
\bibliography{references}

\end{document}
% ============================================================
%  END OF FILE
% ============================================================
